\chapter{Operas}
\label{ch:operas}

\newthought{Use \yamlkey{Operas}} when your calculation can run directly in Python through a
\Operas\ operator---no input files, no shell commands, just a function call. It is the fastest way
to prototype a scan, and the backend the built-in quick-starts use.

\section{When to choose Operas}
\label{sec:op-when}

Pick \yamlkey{Operas} when:

\begin{itemize}
  \item the calculation already exists as a registered operator;
  \item you want a lightweight prototype before building a file-based calculator;
  \item the work is algebraic or array-based and needs no external executable.
\end{itemize}

\noindent Choose \yamlkey{Calculators} (Chapter~\ref{ch:calculators}) instead when you must drive
an external program through files.

\section{The section shape}
\label{sec:op-shape}

\begin{yamlwin}
Operas:
  make_paraller: 4
  Functions:                      # optional: expose operators to expressions
    - name: helper.eggbox2d
      alias: eggbox2d
      inputs: [x, y]
  Modules:
    - name: EggBox
      operator: "helper.eggbox2d"
      call_mode: call
      required_modules: []
      input:
        - {name: x, expression: "xx * Pi"}
        - {name: y, expression: "yy * Pi"}
      kwargs: {}
      output:
        - {name: z, entry: z}
\end{yamlwin}

\begin{description}[style=nextline, leftmargin=1.2em]
  \item[\yamlkey{make\_paraller}] worker count (default 4; keep the legacy spelling).
  \item[\yamlkey{Functions}] an optional whitelist that exposes selected \Operas\ functions to
        symbolic expressions (Section~\ref{sec:op-functions}).
  \item[\yamlkey{Modules}] the operator modules that make up the workflow.
\end{description}

\section{Module keys}
\label{sec:op-keys}

\begin{description}[style=nextline, leftmargin=1.2em]
  \item[\yamlkey{name}] a unique module name, used in logs and dependency resolution.
  \item[\yamlkey{operator}] the registered \Operas\ operator to call, e.g.
        \code{helper.eggbox2d}.
  \item[\yamlkey{call\_mode}] how to invoke it: \code{call} (default) or \code{acall} (async).
  \item[\yamlkey{required\_modules}] upstream modules that must run first (empty list if none).
  \item[\yamlkey{input}] the operator's inputs (see below).
  \item[\yamlkey{kwargs}] extra keyword arguments passed to the operator.
  \item[\yamlkey{output}] the observables to read back, each a \yamlkey{name} and the source
        \yamlkey{entry} in the operator's result.
\end{description}

\paragraph{The \yamlkey{input} list} Each item is either a plain string---``forward this
observable under the same name''---or a mapping with explicit keys:

\begin{description}[style=nextline, leftmargin=1.2em]
  \item[\yamlkey{name}] the operator's argument name;
  \item[\yamlkey{expression}] an optional symbolic expression to compute the value
        (Chapter~\ref{ch:symbols});
  \item[\yamlkey{entry}] an optional source field, to forward an existing observable under a new
        name.
\end{description}

\begin{yamlwin}
input:
  - "mh"                                   # forward observable mh unchanged
  - {name: x, expression: "xx * Pi"}       # compute from current observables
  - {name: sigma, entry: "dd.sigma_si"}    # rename an existing field
\end{yamlwin}

\section{Exposing operator functions to expressions}
\label{sec:op-functions}

The optional \yamlkey{Functions} list makes chosen \Operas\ functions callable inside symbolic
expressions (likelihoods, mappings). Each item supports:

\begin{description}[style=nextline, leftmargin=1.2em]
  \item[\yamlkey{name}] the full registered function name;
  \item[\yamlkey{alias}] an optional short name to use in expressions;
  \item[\yamlkey{inputs}] the ordered input names for the generated wrapper.
\end{description}

\noindent This is handy when you want to call an operator directly from a \yamlkey{LogLikelihood}
term rather than as a workflow module.

\section{A worked example}
\label{sec:op-example}

The quick-start's \code{EggBox} module (Chapter~\ref{ch:quickstart}) is a complete Opera: it runs
the operator \code{helper.eggbox2d} on the scan variables \yamlkey{x} and \yamlkey{y} and reads
the result \yamlkey{z}:

\begin{yamlwin}
Operas:
  make_paraller: 8
  Modules:
    - name: EggBox
      operator: "helper.eggbox2d"
      call_mode: call
      required_modules: []
      input:
        - {name: x, expression: "x"}
        - {name: y, expression: "y"}
      output:
        - {name: z, entry: z}
\end{yamlwin}

\noindent The observable \yamlkey{z} then flows into the likelihood, exactly as a calculator's
output would.

\section{Summary}
\label{sec:op-summary}

\yamlkey{Operas} runs in-process operators: name the \yamlkey{operator}, map its \yamlkey{input}
(strings, expressions, or renames), and read its \yamlkey{output}. It is the quickest workflow
backend and is covered further---including the operator catalog---in
Chapter~\ref{ch:jarvis-operas}.
